
\medskip
\noindent\textbf{0.3.}\footnote{Editors' note: This paragraph has been
modified to take account of the additions made in \SGARefLink{sga3:VIIB:0.2}{0.2}.}
\phantomsection
\label{sga3:VIIB:0.3}
Let \(L\) and \(M\) be two pseudocompact \(A\)-modules. The functor
\[
  \operatorname{LF}(A)\longrightarrow(\mathrm{Ab}),
  \qquad
  N\longmapsto\operatorname{Bil}_{c}(L\times M,N),
\]
where \(\operatorname{Bil}_{c}(L\times M,N)\) denotes the abelian group
of continuous \(A\)-bilinear maps from \(L\times M\) to \(N\), is left
exact.

By \SGARefLink{sga3:VIIB:0.2.E}{0.2.E}, there therefore exists a pseudocompact \(A\)-module
\(L\widehat{\otimes}_{A}M\), unique up to unique isomorphism, representing
this functor; that is, for every object \(N\) of
\(\operatorname{LF}(A)\), there is a functorial isomorphism
\[
  \operatorname{Hom}_{c}
    \bigl(L\widehat{\otimes}_{A}M,N\bigr)
  \simeq
  \operatorname{Bil}_{c}(L\times M,N).
\]

Moreover, \(L\widehat{\otimes}_{A}M\) identifies with the projective
limit
\[
  P=\widehat{L\otimes_{A}M}
\]
of the discrete \(A\)-modules
\[
  (L/L')\otimes_{A}(M/M'),
\]
where \(L'\) and \(M'\) range respectively over the open submodules of
\(L\) and \(M\).

Indeed, let
\(\varphi:L\times M\to N\) be a continuous bilinear map into a discrete
\(A\)-module \(N\) of finite length. By \SGARefLink{sga3:VIIB:0.3.1}{Lemma~0.3.1} below, there exist
open submodules \(L'\) and \(M'\) of \(L\) and \(M\) such that
\[
  \varphi(L'\times M)=\varphi(L\times M')=\{0\}.
\]
This means that the map
\[
  \overline{\varphi}:L\otimes_{A}M\longrightarrow N
\]
induced by \(\varphi\) has the form \(\varphi'\circ p\), where \(p\) is
the canonical projection from \(L\otimes_{A}M\) onto
\((L/L')\otimes_{A}(M/M')\). If \(\widehat{\varphi}\) denotes the
composite
\[
  P
  \longrightarrow
  (L/L')\otimes_{A}(M/M')
  \xrightarrow{\ \varphi'\ }
  N,
\]
then the map \(\varphi\mapsto\widehat{\varphi}\) is a functorial
bijection from \(\operatorname{Bil}_{c}(L\times M,N)\) onto
\(\operatorname{Hom}_{c}(P,N)\). Hence
\[
  P\simeq L\widehat{\otimes}_{A}M.
\]

Thus the pseudocompact module \(L\widehat{\otimes}_{A}M\) is the
separated completion of \(L\otimes_{A}M\) for the linear topology
defined by the kernels of the canonical projections
\[
  L\otimes_{A}M
  \longrightarrow
  (L/L')\otimes_{A}(M/M'),
\]
and it will be called the \emph{completed tensor product} of \(L\) and
\(M\).

If \(x\in L\) and \(y\in M\), the image of \(x\otimes_{A}y\) in
\(L\widehat{\otimes}_{A}M\) will be denoted by
\(x\widehat{\otimes}_{A}y\).

\medskip
\noindent\textbf{Lemma 0.3.1.}
\phantomsection
\label{sga3:VIIB:0.3.1}
\textit{Let \(L\), \(M\), and \(N\) be pseudocompact \(A\)-modules,
with \(N\) of finite length. If
\(\varphi:L\times M\to N\) is a continuous \(A\)-bilinear map, there
exist open submodules \(L'\) and \(M'\) of \(L\) and \(M\) such that}
\[
  \varphi(L'\times M)=\varphi(L\times M')=\{0\}.
\]

\noindent\textit{Proof.}
The inverse image \(\varphi^{-1}(0)\) is an open neighborhood of
\((0,0)\), and therefore contains an open set of the form
\(L_{1}\times M_{1}\), where \(L_{1}\) and \(M_{1}\) are open
submodules of \(L\) and \(M\). Since \(L/L_{1}\) has finite length,
there are elements \(x_{1},\ldots,x_{r}\) of \(L\) such that
\[
  L_{1}+Ax_{1}+\cdots+Ax_{r}=L.
\]
If \(M'\subset M_{1}\) is sufficiently small, then
\(\varphi(x_i,M')=0\) for every \(i\), since the map
\(y\mapsto\varphi(x_i,y)\) is continuous. It follows that
\(\varphi(L,M')=\{0\}\). Likewise,
\(\varphi(L',M)=\{0\}\) when \(L'\) is sufficiently small.
\hfill\(\square\)

\medskip
\noindent\textbf{Corollary 0.3.1.1.}\footnote{Editors' note: This
corollary has been added.}
\phantomsection
\label{sga3:VIIB:0.3.1.1}
\textit{Let \(M\) be a pseudocompact \(A\)-module.}
\begin{enumerate}
  \item[\textup{(i)}]
  \textit{For every open submodule \(M'\), there exists an open ideal
  \(\mathfrak l\) of \(A\) such that}
  \[
    \overline{\mathfrak lM}\subset M'.
  \]

  \item[\textup{(ii)}]
  \textit{Consequently,}
  \[
    M\simeq
    \varprojlim_{\mathfrak l}
      M/\overline{\mathfrak lM},
  \]
  \textit{where \(\mathfrak l\) ranges over the filtered projective
  system of open ideals of \(A\), and every
  \(M/\overline{\mathfrak lM}\) is endowed with the quotient topology,
  which makes it a pseudocompact module; cf.\ \SGARefLink{sga3:VIIB:0.2.D}{0.2.D}.}
\end{enumerate}

\noindent\textit{Proof.}
Consider the map
\[
  \varphi:A\times M\longrightarrow M/M',
  \qquad
  (a,m)\longmapsto am+M'.
\]
By \SGARefLink{sga3:VIIB:0.3.1}{0.3.1}, there is an open ideal \(\mathfrak l\) of \(A\) such that
\(\mathfrak lM\subset M'\). Since \(M'\) is also closed, it contains
\(\overline{\mathfrak lM}\). Since the intersection of the open
submodules of \(M\) is zero, it follows that
\[
  \bigcap_{\mathfrak l}\overline{\mathfrak lM}=(0).
\]
On the other hand, by \SGARefLink{sga3:VIIB:0.2.D}{0.2.D} the map
\[
  \phi:M\longrightarrow
  \varprojlim_{\mathfrak l}M/\overline{\mathfrak lM}
\]
is surjective. By the proof of \SGARefLink{sga3:vii:proposition:0-2-b}{0.2.B}, \(\phi\) therefore induces an
isomorphism
\[
  M/\operatorname{Ker}(\phi)
  \xrightarrow{\ \sim\ }
  \varprojlim_{\mathfrak l}M/\overline{\mathfrak lM}.
\]
But we have just seen that
\[
  \operatorname{Ker}(\phi)
  =
  \bigcap_{\mathfrak l}\overline{\mathfrak lM}
  =
  (0).
\]
\hfill\(\square\)

\medskip
\noindent\textbf{Remark 0.3.1.2.}\footnote{Editors' note: This remark
has been added.}
\phantomsection
\label{sga3:VIIB:0.3.1.2}
The completed tensor product satisfies the usual associativity
condition: if \(L\), \(M\), and \(P\) are pseudocompact
\(A\)-modules, there is a canonical isomorphism
\[
  (L\widehat{\otimes}_{A}M)\widehat{\otimes}_{A}P
  \simeq
  L\widehat{\otimes}_{A}(M\widehat{\otimes}_{A}P).
\]
Indeed, both objects represent the functor which associates to every
object \(N\) of \(\operatorname{LF}(A)\) the abelian group of
continuous \(A\)-trilinear maps from \(L\times M\times P\) to \(N\).

\medskip
\noindent\textbf{0.3.2.}
\phantomsection
\label{sga3:VIIB:0.3.2}
Let
\[
  L'\xrightarrow{\,f\,}L\xrightarrow{\,g\,}L''\longrightarrow0
\]
be an exact sequence, and let \(M\) be an object of
\(\operatorname{PC}(A)\). It is clear that for every object \(N\) of
\(\operatorname{LF}(A)\), the induced sequences
\[
\begin{array}{ccccccc}
0 &\longrightarrow&
  \operatorname{Bil}_{c}(L''\times M,N)
&\longrightarrow&
  \operatorname{Bil}_{c}(L\times M,N)
&\longrightarrow&
  \operatorname{Bil}_{c}(L'\times M,N)
\\[0.35em]
&&\Big\Vert&&\Big\Vert&&\Big\Vert
\\[-0.1em]
0 &\longrightarrow&
  \operatorname{Hom}_{c}(L''\widehat{\otimes}_{A}M,N)
&\longrightarrow&
  \operatorname{Hom}_{c}(L\widehat{\otimes}_{A}M,N)
&\longrightarrow&
  \operatorname{Hom}_{c}(L'\widehat{\otimes}_{A}M,N)
\end{array}
\]
are exact. By \SGARefLink{sga3:VIIB:0.2.E}{0.2.E}, this is equivalent to saying that the sequence
\begin{equation}
  L'\widehat{\otimes}_{A}M
  \xrightarrow{\,f\widehat{\otimes}M\,}
  L\widehat{\otimes}_{A}M
  \xrightarrow{\,g\widehat{\otimes}M\,}
  L''\widehat{\otimes}_{A}M
  \longrightarrow0
  \tag{\(*\)}
  \label{sga3:VIIB:0.3.2.star}
\end{equation}
is exact. Hence:

\medskip
\noindent\textbf{Corollary.}
\phantomsection
\label{sga3:VIIB:0.3.2.cor}
\textit{For every pseudocompact \(A\)-module \(M\), the functor
\(-\widehat{\otimes}_{A}M\) is right exact.}

In particular, take \(L=A\), let \(f\) be the inclusion of a closed
ideal \(\mathfrak a\) into \(A\), and let \(g\) be the canonical
projection of \(A\) onto \(A/\mathfrak a\). We may then identify
\(A\widehat{\otimes}_{A}M\) with \(M\) by means of the map
\[
  x\widehat{\otimes}_{A}m\longmapsto xm.
\]
Since the image of
\(\mathfrak a\widehat{\otimes}_{A}M\) is closed in \(M\)
(cf.\ \SGARefLink{sga3:vii:proposition:0-2-b}{0.2.B}), and the image of
\(\mathfrak a\otimes_{A}M\) is dense in
\(\mathfrak a\widehat{\otimes}_{A}M\), the image of
\(f\widehat{\otimes}M\) is precisely the closure
\(\overline{\mathfrak aM}\) of \(\mathfrak aM\) in \(M\).
The exact sequence \(\textup{(\(*\))}\) therefore gives the
isomorphism
\[
  (A/\mathfrak a)\widehat{\otimes}_{A}M
  \xrightarrow{\ \sim\ }
  M/\overline{\mathfrak aM}.
\]

\medskip
\noindent\textbf{0.3.3. Nakayama's Lemma.}
\phantomsection
\label{sga3:VIIB:0.3.3}
\textit{Let \(A\) be a pseudocompact ring, \(M\) a pseudocompact
\(A\)-module, and \(\mathfrak a\) an ideal of \(A\) contained in the
radical \(r(A)\). Then}
\[
  \overline{\mathfrak aM}=M
  \qquad\Longrightarrow\qquad
  M=0.
\]

\noindent\textit{Proof.}
Suppose that
\(\overline{\mathfrak aM}=M\).
Let \(M'\) be an open submodule of \(M\), and put \(M''=M/M'\).
Since \(M''\) is discrete, \(\mathfrak aM''\) is closed in \(M''\),
and hence equals \(\overline{\mathfrak aM''}\). By \SGARefLink{sga3:VIIB:0.3.2}{0.3.2}, the
canonical map
\[
  M/\overline{\mathfrak aM}
  \longrightarrow
  M''/\overline{\mathfrak aM''}
\]
is surjective, so
\[
  \mathfrak aM''
  =
  \overline{\mathfrak aM''}
  =
  M''.
\]
Since \(M''\) is a finitely generated \(A\)-module and
\(\mathfrak a\subset r(A)\), the usual Nakayama lemma implies that
\(M''=0\). Consequently, every open submodule \(M'\) of \(M\) is
equal to \(M\), and therefore \(M=0\).
\hfill\(\square\)

\medskip
\noindent\textbf{0.3.4.}
\phantomsection
\label{sga3:VIIB:0.3.4}
The usual consequences of Nakayama's lemma follow.

\medskip
\noindent\textbf{Corollary.}
\phantomsection
\label{sga3:VIIB:0.3.4.cor}
\textit{Let \(\mathfrak a\) be a closed ideal contained in \(r(A)\),
and let \(f:M\to N\) be a morphism of pseudocompact \(A\)-modules.}
\begin{enumerate}
  \item[\textup{(i)}]
  \textit{The map \(f\) is surjective if the induced map}
  \[
    f':
    M/\overline{\mathfrak aM}
    \longrightarrow
    N/\overline{\mathfrak aN}
  \]
  \textit{is surjective.}

  \item[\textup{(ii)}]
  \textit{If \(N\) is projective, then \(f\) is an isomorphism whenever
  \(f'\) is an isomorphism.}
\end{enumerate}

\noindent\textit{Proof.}
Part \textup{\SGARefLink{sga3:vii:paragraph:5-1:equation:i}{(i)}} follows by applying \SGARefLink{sga3:VIIB:0.3.3}{Lemma~0.3.3} to
\(\operatorname{Coker}(f)\). For \textup{\SGARefLink{sga3:vii:paragraph:5-1:equation:ii}{(ii)}}, suppose that \(f'\) is
an isomorphism. By \textup{\SGARefLink{sga3:vii:paragraph:5-1:equation:i}{(i)}}, \(f\) is then surjective, and hence
has a section. Apply \SGARefLink{sga3:VIIB:0.3.3}{Lemma~0.3.3} to \(\operatorname{Ker}(f)\).
\hfill\(\square\)\par\smallskip

When \(A\) is local with maximal ideal \(\mathfrak m\), one may also
deduce from \SGARefLink{sga3:VIIB:0.3.3}{0.3.3} the following exchange theorem.

\medskip
\noindent\textbf{Theorem.}
\phantomsection
\label{sga3:VIIB:0.3.4.exchange}
\textit{Let \(A\) be a local pseudocompact ring, let \(\mathfrak m\)
be its maximal ideal, let \(M\) be a topologically free \(A\)-module
with pseudobasis \((m_i)_{i\in I}\) \textup{(\SGARefLink{sga3:VIIB:0.2.1}{0.2.1})}, and let \(N\)
be a closed direct summand of \(M\). There is a pseudobasis of \(M\)
consisting of elements of \(N\) together with some of the \(m_i\).}

\noindent\textit{Proof.}
This is clear when \(A\) is a field: use the duality of \SGARefLink{sga3:VIIB:0.2.2}{0.2.2} and apply
the usual exchange theorem. Consequently,
\[
  N/\overline{\mathfrak mN}
\]
has a complementary topologically free \(A/\mathfrak m\)-module with
pseudobasis \((\overline{m}_i)_{i\in J}\), where
\(\overline{m}_i\) is the image of \(m_i\) in
\[
  M/\overline{\mathfrak mM},
\]
and \(J\) is a subset of \(I\).
If
\[
  P=\prod_{i\in J}Am_i,
\]
then the canonical map
\[
  N\oplus P\longrightarrow M
\]
is ``bijective modulo \(\mathfrak m\).'' It is therefore bijective by
what precedes. For another proof, see \textup{[CA]}, \S~IV.2,
Proposition~8.
\hfill\(\square\)

